%% sections/introduction.tex
\section{Introduction}
\label{sec:intro}

Experimental asset markets in the tradition of
\citet{smith1988bubbles} have produced one of the most replicated
departures from rational-expectations pricing in all of economics:
even when subjects know, in closed form, that the fundamental value
of a finite-lived asset declines deterministically with the
remaining dividend stream, transaction prices systematically
overshoot the fundamental before collapsing near the terminal
period. \citet{dufwenberg2005bubbles}---hereafter
DLM~2005---established a by now canonical result in this literature:
when two-thirds or more of the subject pool consists of traders who
have played the game before, the bubble~disappears. Experience, rather
than cognition or incentives, is what eventually kills the bubble.

A parallel and much more recent literature replaces the human
subjects with AI agents. \citet{lopezlira2025llm} proposes an
expected-utility framework in which each AI-agent trader carries a
persistent valuation bias $b_i$, a risk preference over wealth
$U\in\{\Uloving,\Uneutral,\Uaverse\}$, and a belief-blending channel
that lets agents exchange private valuations and update a
trust score for each sender by exponential moving average. Lopez-Lira's
agents optimise expected utility over a small action set---hold,
cross the bid, cross the ask, post passively---and the resulting
population can be dialled between skeptical and naive, truthful and
deceptive, risk-averse and risk-loving by simple categorical knobs.
The Lopez-Lira design is a natural vehicle for studying~how a
heterogeneous AI-agent population prices a DLM asset, but it leaves
one channel unspecified: \emph{how should each utility agent
translate its prior and the messages it has received into a
posterior valuation that it then carries into the next period?}
Every downstream instrument---the expected-utility score, the
order chosen, the trust EMA, the order-matching step, the dividend
credit---consumes that posterior~$V_i^{\text{post}}$ and is
otherwise unaffected by the belief-update channel.

This paper contributes a computational testbed that holds the
market substrate, the agent population, the starting prior, and the
action-choice rule strictly fixed, and compares three distinct
\emph{plans} for the per-period belief update.

\paragraph{Plan I (algorithmic baseline).} The posterior is a
convex combination of the agent's own prior and the arithmetic
mean of foreign peer claims received during the period, with a
weight $w$ that ramps endogenously with experience,
\begin{equation}
V_i^{\text{post}} \;=\; w\cdot V_i^{\text{prior}}
\;+\;(1-w)\cdot\frac{1}{|\mathcal{M}_i|}\sum_{m\in\mathcal{M}_i}\hat v_m,
\qquad w = 0.6 + 0.1\,\min(3,\,k_i),
\label{eq:plan1}
\end{equation}
where $k_i$ is the agent's \texttt{roundsPlayed} counter. A novice
($k_i=0$) moves appreciably toward peers ($w=0.60$), a
three-round veteran ($k_i\ge3$) barely moves ($w=0.90$). No network
calls are placed; the update is deterministic under the seeded
PRNG.

\paragraph{Plan II (LLM update with explicit utility forms).}
At every period boundary the engine fires one parallel
\texttt{chat/completions} call per utility agent. The user prompt
encodes the market constants, the agent's risk preference and bias
mode, the peer claims accumulated during the period, the agent's
mark-to-market wealth, and---critically---the \emph{closed-form}
expression of the agent's risk-utility functional:
\begin{equation}
\pi_i^{\text{II}} \;\supset\; \bigl\{\,U_L(w)=(w/w_0)^2,\;U_N(w)=w/w_0,\;U_A(w)=\sqrt{w/w_0}\,\bigr\}.
\label{eq:plan2-prompt}
\end{equation}
The model is asked to return a single number in experimental cents,
which the engine parses, clamps to a sanity interval, and writes to
the agent's next-period prior slot. The call is fire-and-forget: if
the network errors or the parse fails, the agent silently degrades
to Plan~I for that period, so reproducibility is preserved even
when the key is absent.

\paragraph{Plan III (LLM update with risk label only).} Identical
wiring to Plan~II, but the user prompt strips the closed-form
expressions out of equation~\eqref{eq:plan2-prompt} and names only
the risk-preference category:
\begin{equation}
\pi_i^{\text{III}} \;\supset\;
\{\text{risk-loving}\,|\,\text{risk-neutral}\,|\,\text{risk-averse}\}.
\label{eq:plan3-prompt}
\end{equation}
Plan~III therefore isolates a narrow, purely linguistic question:
does handing the model an explicit functional form for utility
produce different market dynamics than naming the risk type alone?
The LLM in Plan~III must translate ``risk-averse'' into an implicit
curvature internally; the LLM in Plan~II is handed the curvature
as a formula. Any divergence between the two plans' market statistics
quantifies the prompt-sensitivity of the model's internal
expected-utility reasoning.

\paragraph{Research question.} With the market substrate, the agent
population, the Lopez-Lira prior, the expected-utility action-choice
rule, the order-matching engine, and the seed all held fixed, the
question the testbed is built to answer~is
\begin{equation*}
\text{Plan~I} \;\overset{?}{\approx}\; \text{Plan~II} \;\overset{?}{\approx}\; \text{Plan~III}
\quad\text{in classical bubble statistics and posterior trajectory?}
\end{equation*}
A null result (all three plans produce indistinguishable market
dynamics) would imply that the specific belief-update channel is
absorbed by the downstream expected-utility score and does not
propagate into observable market behaviour. A divergence between
Plan~I and the two LLM plans would implicate the LLM itself as a
source of systematic behavioural difference. A divergence
\emph{between} Plan~II and Plan~III would implicate prompt
structure---specifically, whether handing the model a closed-form
utility function alters its valuation choice---as a source of
behavioural difference that goes beyond the agent's declared risk
type.

Section~\ref{sec:related} places the simulator in the literature on
experimental asset markets, expected-utility AI-agent traders, and
LLM prompt sensitivity. Section~\ref{sec:design} documents the
market structure and the sampling stage. Section~\ref{sec:agents}
gives the precise decision rule for each agent type.
Section~\ref{sec:anchor} presents the three belief-update plans:
the algorithmic blend, the two LLM calls, the prompt templates,
the clamp, the consume-and-cache path, and the
fallback-to-Plan~I~invariant. Section~\ref{sec:results} reports
preliminary findings, and Section~\ref{sec:conclusion} concludes.
